\documentclass[10pt,a4paper]{amsart}
\usepackage[margin=19mm]{geometry}
\usepackage{amsmath,amssymb,mathtools}
\usepackage[scr=rsfs,cal=euler]{mathalfa}
\usepackage{xfrac}
\usepackage[unicode,hidelinks,bookmarks=false]{hyperref}
\newcommand{\ie}{i.e.\ }
\newcommand{\cf}{cf.\ }
\newcommand{\resp}{respectively\ }
\newcommand{\Subsection}{\subsection}
\providecommand{\nobreakdash}{\nobreak\hbox{-}}
\setlength{\emergencystretch}{2em}
\multlinegap0pt
\usepackage{bm}
\usepackage{bbm}
\usepackage{dsfont}
\usepackage{xspace}
\usepackage{cancel}
\usepackage{mathtools}
\usepackage{amsthm}
\makeatletter
\@ifundefined{theorem}{\newtheorem{theorem}{Theorem}}{}
\@ifundefined{lemma}{\newtheorem{lemma}[theorem]{Lemma}}{}
\makeatother
\newcommand{\Z}{\mathbb{Z}}
\title[Positive Lower Density for Hofstadter's $ab-1$ Problem]{Positive Lower Density for Hofstadter's $ab-1$ Problem}
\author{Samuel Korsky}
\date{Source: arXiv:2608.07910v1 (CC BY 4.0)}
\begin{document}
\maketitle

\noindent\emph{Source and licence.} Positive Lower Density for Hofstadter's $ab-1$ Problem. Version arXiv:2608.07910v1 (CC BY 4.0), distributed under the Creative Commons Attribution 4.0 International licence, as recorded in the arXiv licence metadata for this exact version and shown on the abstract page https://arxiv.org/abs/2608.07910v1. Full rights notice: licenses/arxiv-2608.07910-CC-BY-4.0.txt.\par\smallskip

\noindent\textit{Editorial scope.} Counted unit: the Lemma whose source label is \texttt{lem:recurrence} in the article (source0.tex lines 204--206, source label \texttt{lem:recurrence}), delivered together with the article's own complete original proof of it (source0.tex lines 208--282, one \texttt{proof} environment of 75 lines). No mathematics is written, repaired or re-derived: statement and proof are the article's own text line for line. Required context retained uncounted: none. Every definition, display and label of those lines is retained unchanged. Omitted from this package and not redistributed: 1--203, 207--207, 283--497. Nothing inside a retained excerpt is omitted or altered: every retained body is delivered exactly as the article's lines are written, so no omission receipt of any kind accompanies this package. Citations: the retained excerpts carry no citation command at all, so no bibliography entry of the article is transcribed and no reference is redistributed. Reference adaptation: none was needed and none was made; every reference inside the delivered unit resolves inside it, so no unresolved reference or citation is produced. Printed slips kept literal: no printed word, symbol, hypothesis or number of the counted statement or of its proof was altered. Layout and class: the article's own class is \texttt{article} with packages including geometry, amsmath, amssymb, amsthm, booktabs, microtype, xcolor, fontenc, lmodern, hyperref; the wrapper is the lane's pinned \texttt{amsart} editorial wrapper at 10pt on A4, which restates the theorem-like environments the retained text needs and adds the article's own macro definitions that the retained text needs (\texttt{\textbackslash Z}), with extra wrapper packages (bm, bbm, dsfont, xspace, cancel, mathtools). The delivered numbering is the unit's own. Assets: no retained span contains a figure environment, a graphics-inclusion command, a PDF-page inclusion, a table or a TikZ picture; no image file is copied into this package, the package declares no asset dependency and no image file is redistributed with it. Licence: version arXiv:2608.07910v1 of this article is distributed under the Creative Commons Attribution 4.0 International licence, as recorded in the arXiv licence metadata for this exact version and shown on the abstract page https://arxiv.org/abs/2608.07910v1. A full rights notice is delivered at licenses/arxiv-2608.07910-CC-BY-4.0.txt. Characters: every retained excerpt is pure ASCII, so the delivered engine, XeLaTeX with its default Latin Modern text font, reports no missing character.
\begin{lemma}\label{lem:recurrence}
There is a deterministic way to choose among $\alpha_0,\ldots,\alpha_3$ such that the interval state together with the imbalance vector has a positive recurrent state.
\end{lemma}

\begin{proof}
For fixed $\alpha$, let $\xi_\alpha(i)$ be the increment of $H_t$ on leaving $K_i$. Since $P_\alpha$ is irreducible on a finite state space and $\xi_\alpha-d_\alpha$ has stationary mean zero, the Markov-chain Poisson equation
\[
g_\alpha-P_\alpha g_\alpha=\xi_\alpha-d_\alpha
\]
has a bounded solution, coordinatewise. Hence there is a constant $C_0$, independent of $\alpha$, the initial state $i$, and $N\geq1$, such that
\begin{equation}\label{eq:averaging}
\left|
\mathbb E_i\left[\sum_{t=0}^{N-1}\xi_\alpha(I_t)\right]-Nd_\alpha
\right|\leq C_0.
\end{equation}

Since the origin lies in the interior of the convex hull of the four drift vectors, there is $\delta>0$ such that for every unit vector $v\in\mathbb R^3$,
\[
\min_\alpha\langle v,d_\alpha\rangle\leq-\delta.
\]
At each time $kN$, choose $\alpha_k\in\{\alpha_0,\alpha_1,\alpha_2,\alpha_3\}$ and use $P_{\alpha_k}$ for the next $N$ steps. If $H_{kN}=h\neq0$, choose $\alpha_k$ to minimize
\[
\left\langle \frac{h}{|h|},d_\alpha\right\rangle.
\]
Fix a deterministic tie-breaking convention, and take $\alpha_k=\alpha_0$ when $h=0$. Then
\[
X_k=(I_{kN},H_{kN})
\]
is a time-homogeneous Markov chain.

Put
\[
M_0=\max_{\alpha,i}|\xi_\alpha(i)|.
\]
Combining \eqref{eq:averaging} with
\[
|h+z|-|h|\leq \left\langle\frac h{|h|},z\right\rangle+\frac{|z|^2}{2|h|}
\]
gives
\[
\mathbb E\left[
|H_{(k+1)N}|-|H_{kN}|
\,\middle|\,
X_k=(i,h)
\right]
\leq-N\delta+C_0+\frac{M_0^2N^2}{2|h|}.
\]
First choose $N$ with $N\delta>C_0+2$, and then choose $R$ so that the right-hand side is at most $-1$ whenever $|h|>R$.

Let
\[
\mathcal C=\{(i,h):1\leq i\leq19,\ h\in\Z^3,\ |h|\leq R\}
\]
and define
\[
\tau_{\mathcal C}=\inf\{k\geq0:X_k\in\mathcal C\},
\qquad
\tau_{\mathcal C}^+=\inf\{k\geq1:X_k\in\mathcal C\}.
\]
The drift estimate and optional stopping give
\[
\mathbb E_x\left[\tau_{\mathcal C}\right]\leq |h|
\]
for $x=(i,h)\notin\mathcal C$. If $x\in\mathcal C$, then $|H_N|\leq R+M_0N$, and the Markov property gives
\[
\begin{aligned}
\mathbb E_x\left[\tau_{\mathcal C}^+\right]
&=1+\mathbb E_x\left[
\mathbf 1_{\{X_1\notin\mathcal C\}}
\mathbb E_{X_1}\left[\tau_{\mathcal C}\right]
\right]\\
&\leq1+\mathbb E_x\left[
\mathbf 1_{\{X_1\notin\mathcal C\}}|H_N|
\right]\\
&\leq1+R+M_0N.
\end{aligned}
\]
The chain induced by successive visits to $\mathcal C$ has finite state space. Choose $s$ in one of its recurrent classes. Its return time in the induced chain has finite mean, and the uniform bound above on the expected time between successive visits to $\mathcal C$ implies that its return time in $(X_k)$ also has finite mean. Thus $s$ is positive recurrent.
\end{proof}

\end{document}
